\begin{table}[t]\centering\footnotesize
\caption{Incremental computation costs (seconds) under the shared
continuation protocol. Training-label and common-validation costs are
measured separately on the same CPU protocol; training times are
archived continuation medians over the weight grid and seeds, and
explanation times are sequential medians for the selected models.}
\label{tab:learning-costs}
\begin{tabular}{lrrrr}\toprule Labels & Training labels & Common validation & Continuation & Explanation\\\midrule
Flow only & 0.00 & 6.45 & 246.0 & 7.34\\
Response family & 16.12 & 6.45 & 155--425 & 7.13\\
Endpoint (edge, cut) & 25.75 & 6.45 & 441--497 & 7.27--7.59\\
\bottomrule\end{tabular}\end{table}
Training labels comprise 768 endpoint states for either endpoint
objective, or 384 center states and their sensitivities for the three
response-family objectives (unweighted, weighted, boundary-seminorm),
which share one label set. Every method additionally uses the same 192
endpoint states for validation; unused center-validation derivatives
are excluded. Existing flow data, pretrained models and the common test
reference are excluded from these incremental costs. The 984 test
states and sensitivities took 74.4\,s in the archived reference run.
Explanation timing covers all 24 paths, forward evaluation, six
directional derivatives per state, quadrature and node aggregation,
excluding model loading. Archived continuation times include
initialization, validation and saving under concurrent jobs, pooled
across the original run and the corrected-weight re-run under lighter
load; they describe the actual training expense and do not establish a
method speed advantage.
